\section{Generalized-circle blocks}

\begin{definition}
A \emph{generalized-circle block} of $S$ is a maximal subset of at least three selected points lying either on one line or on one proper circle.  Write
\[
 q=\#\{\text{proper-circle blocks}\},\qquad
 L=\#\{\text{rich-line blocks}\},\qquad B=q+L,
\]
and let $b_s$ be the number of generalized-circle blocks of size $s$.
\end{definition}

The key bookkeeping fact is immediate but powerful.

\begin{lemma}[Exact triple partition]\label{lem:triplepartition}
Every three-element subset of $S$ belongs to exactly one maximal generalized-circle block.  Consequently
\begin{equation}\label{eq:triplepartition}
 \sum_{s\ge3}\binom{s}{3}b_s=\binom n3.
\end{equation}
\end{lemma}

\begin{proof}
Three noncollinear points determine a unique proper circle, whereas three collinear points determine a unique line.  Maximal extension is unique.  Indeed, two distinct lines meet in at most one point, two distinct proper circles meet in at most two points, and a line and a proper circle meet in at most two points.
\end{proof}

Two further pair-budget observations recur throughout the finite arguments.  Distinct rich lines use disjoint unordered point pairs, and distinct generalized-circle blocks meet in at most two selected points.  These yield clique, degree and convexity obstructions after a block profile is fixed.

